% supp_datasets.tex -- Supplementary Section S6: the outbreak records (body only; \input by supplement.tex).
% Paths in "% src:" comments are relative to the repository root.
% Tables: the three table bodies between the markers ">>> GENERATED" / "<<< GENERATED" are written by
%   code/appB_dataset_table.py (from data/bsc_outbreaks/outbreaks.csv and data/bsc_outbreaks/tables/*.csv);
%   all other numbers of this appendix are in data/appB_dataset_numbers.json (same script).
% Figure: code/appB_dataset_fig_attack.py -> figures/appB_dataset_attack_{en,zh}.pdf.
\section{The outbreak records}\label{appB_dataset:sec}

This section documents the two datasets of outbreak records used in Section~\ref{M-s8_outbreaks:sec} of the main
text. For the first dataset, 27 records of outbreaks and exposure cohorts (Section~\ref{appB_dataset:sec:first}), it
says how they were found, what was recorded, how reliable each record is, and which numbers were digitised or
reconstructed. The second dataset, eleven records of ten outbreaks with positions, follows in
Section~\ref{appB_dataset:sec:second}.

\subsection{The first dataset: 27 records of outbreaks and exposure cohorts}\label{appB_dataset:sec:first}
% src: data/bsc_outbreaks/outbreaks.csv (27 rows); data/appB_dataset_numbers.json -> n_records

\paragraph{Search and inclusion.}
Candidate outbreaks were located by web search, with one set of queries per scene (office and call centre;
supermarket and retail; classroom and school; metro, bus, train and aircraft; restaurant; choir), and by
following the citations of the papers found. This is not a systematic review and the coverage is not
exhaustive. A record was included if (a) it is described in a peer-reviewed article or an official
surveillance report, (b) the DOI of that source is registered, and (c) the source gives at least one of: the
number of exposed people, positions at seat or zone level, the geometry of the room, the duration of exposure.
% src: code/bsc_outbreaks/scripts/build_dataset.py; data/bsc_outbreaks/sources.md
The result is 27 records from 31 sources: 5 offices and workplaces, 4 retail settings, 5 classrooms and
class-like rooms, 7 vehicle cabins, 2 restaurants, 2 choir or church events, and 2 records of large settings that
lie outside the scope of the model (a cruise ship and worker dormitories). These two are recorded for reference
only, as examples of large shared settings with long exposure that a model of one closed room does not cover, and
take part in no comparison.
% src: data/appB_dataset_numbers.json -> records_by_group, sources.n
Three records without a denominator (S3, C4, T7) and one with a reported attack rate but no transcribed
counts (T6: 54 persons aboard, 72\pct{} of the passengers ill, from the abstract) were kept as context for their
settings; like the two out-of-scope records, they take no part in any quantitative comparison.
% src: data/appB_dataset_numbers.json -> records_without_denominator (S3, C4, T7), records_with_reported_rate_but_no_counts (T6), records_with_counts (23); data/bsc_outbreaks/outbreaks.csv -> T6 (occupancy_n 54, attack_rate_reported); data/refs_search.json (manual_checks.moser1979)
No documented outbreak in a metro carriage with an exposed denominator was found; buses, a coach, a minibus,
a train cohort and one flight stand in for that scene. The supermarket has one record with denominators for
staff and customers (S1) and none with geometry, positions or dwell times.
% src: data/bsc_outbreaks/outbreaks.csv (scene groups); notes/VERIFICATION.md (Section 4.4, outbreak dataset, item 3: unverifiable)

\paragraph{What was recorded.}
Each record has 51 fields: setting, place, dates and pathogen; the sources and how each was read; type and
duration of exposure; room geometry, occupancy, ventilation and mask use; the index case; the numbers
exposed and infected, with an alternative count where sources or case definitions differ; the attack rate
with its exact (Clopper--Pearson) 95\,\% interval; the spatial pattern reported; a quality grade; a list
of the fields that were derived and not transcribed; and the mapping to the model.
% src: data/bsc_outbreaks/outbreaks.csv (header, 51 columns); data/appB_dataset_numbers.json -> n_columns
Table~\ref{appB_dataset:tab:records} lists the records and Fig.~\ref{appB_dataset:fig:attack} shows the
attack rates. Three records concern other pathogens (O5 tuberculosis, C4 measles, T6 influenza).
% src: data/appB_dataset_numbers.json -> other_pathogens

{\footnotesize
\setlength{\tabcolsep}{3pt}
\setlength{\LTcapwidth}{\textwidth}
\begin{longtable}{@{}p{0.45cm}>{\raggedright}p{3.2cm}>{\raggedright}p{2.25cm}>{\raggedright}p{2.3cm}>{\raggedright}p{1.75cm}>{\raggedright}p{2.5cm}p{0.45cm}>{\raggedright\arraybackslash}p{1.4cm}@{}}
\caption{The 27 records of the outbreak dataset. Counts exclude the index case unless marked $\dagger$.
Attack rates and exact (Clopper--Pearson) 95\,\% intervals are computed from the counts in the third column;
`(or \dots)' gives the count of a second source or case definition. Exposure is the duration of the event,
or of each exposure where it was repeated. Grades are those assigned when the dataset was built (criteria
in the text). Record T4 pools 2,334 index patients.}
\label{appB_dataset:tab:records}\\
\toprule
ID & Setting (place) & Cases / exposed & Attack rate, \% (95\,\% interval) & Exposure & Ventilation & Gr. & Sources \\
\midrule
\endfirsthead
\multicolumn{8}{@{}l}{\emph{Table~\thetable{} (continued)}}\\
\toprule
ID & Setting (place) & Cases / exposed & Attack rate, \% (95\,\% interval) & Exposure & Ventilation & Gr. & Sources \\
\midrule
\endhead
\midrule
\multicolumn{8}{r@{}}{\emph{continued on the next page}}\\
\endfoot
\bottomrule
\addlinespace[3pt]
\multicolumn{8}{@{}p{15.6cm}@{}}{$\dagger$ numerator and denominator include the index or primary case(s).
$\ddagger$ upper bound on the number infected in the restaurant (see text).
$*$ graded A, but not all grade-A criteria are met (see text).
$\S$ numerator derived from the reported attack rate (26.3\pct{} of 217); the record pools the classes of eight
infected instructors in 12 facilities and is not a single-index event.
Reading access to the source: $^{\mathrm{a}}$ abstract only; $^{\mathrm{e}}$ publisher's page available only as an automated text extraction; $^{\mathrm{m}}$ accepted manuscript; unmarked, full text.}\\
\endlastfoot
% src: data/bsc_outbreaks/outbreaks.csv; intervals recomputed and checked against the columns ci95_low, ci95_high
% >>> GENERATED records (code/appB_dataset_table.py; do not edit by hand)
\addlinespace[2pt]\multicolumn{8}{@{}l}{\emph{Offices and workplaces}}\\
O1 & Call centre, 11th floor, open-plan desks (Seoul) & 94/216$^{\dagger}$ & 43.5 (36.8--50.4) & working days over about 2 weeks & not reported & B & \citep{park2020} \\
O2 & Team of 13 in an open-plan office (Zurich) & 8/12 (or 10/12) & 66.7 (34.9--90.1) & 11 h & mechanical, about 1 air change per hour & B & \citep{weissberg2020} \\
O3 & Public-facing office, three open-plan floors (England) & 22/40$^{\dagger}$ & 55.0 (38.5--70.7) & weeks & CO$_2$ typically $\le$1000 ppm & C & \citep{nicholls2024}$^{\mathrm{a}}$ \\
O4 & Beef-processing line, fixed work stations (Germany) & 20/78 & 25.6 (16.4--36.8) & 3 shifts & $<$1 air change per hour, recirculating cooling fans & B & \citep{guenther2020} \\
O5 & Office building, tuberculosis (Massachusetts) & 27/67 & 40.3 (28.5--53.0) & 160 h (4 weeks) & about 7 l/s outdoor air per occupant & C & \citep{nardell1991}$^{\mathrm{a}}$ \\
\addlinespace[2pt]\multicolumn{8}{@{}l}{\emph{Supermarket and retail}}\\
S1 & Supermarket, staff (Liaocheng) & 11/120$^{\dagger}$ & 9.2 (4.7--15.8) & not reported & not reported & B & \citep{tian2021} \\
S2 & Department store, sales staff (Tianjin) & 6/194$^{\dagger}$ (or 7/194) & 3.1 (1.1--6.6) & not reported & described as poorly ventilated & C & \citep{chen2020tianjin,li2022twoclusters}, \citep{jiang2021deptstore}$^{\mathrm{a}}$ \\
S3 & Shopping mall, 8 floors (Wenzhou) & no denominator & -- & not reported & not reported & C & \citep{cai2020} \\
S4 & Grocery store, staff, PCR screen (Massachusetts) & 21/104$^{\dagger}$ & 20.2 (13.0--29.2) & not reported & not reported & C & \citep{lan2021}$^{\mathrm{a}}$ \\
\addlinespace[2pt]\multicolumn{8}{@{}l}{\emph{Classrooms and class-like rooms}}\\
C1 & Primary-school classroom, 24 pupils in 5 rows (Marin County) & 12/24 & 50.0 (29.1--70.9) & school days; hours not reported & portable HEPA filter; doors and windows open & A$^{*}$ & \citep{lamhine2021} \\
C2 & High school, 35--38 pupils per classroom of 39--49 m$^2$ (Jerusalem) & 153/1{,}161$^{\dagger}$ & 13.2 (11.3--15.3) & 8--9 school days of 6.3--6.7 h & split air conditioning, recirculating & B & \citep{steinzamir2020} \\
C3 & 20 primary schools, masks worn, seats 0.9 m apart (Utah) & 5/728 & 0.69 (0.22--1.60) & school days; not quantified & not quantified & B & \citep{hershow2021} \\
C4 & Primary school, 14 classrooms on one recirculating system, measles (New York State) & 28 cases; no denominator & -- & not in abstract & largely recirculated air & C & \citep{riley1978}$^{\mathrm{a}}$ \\
C5 & Fitness dance classes in a 60 m$^2$ room (Cheonan) & 57/217$^{\S}$ & 26.3 (20.5--32.7) & 50 min per class & not reported & B & \citep{jang2020} \\
\addlinespace[2pt]\multicolumn{8}{@{}l}{\emph{Vehicle cabins (stand-ins for the metro carriage)}}\\
T1 & Bus, 15 rows of seats (Zhejiang) & 23/67 & 34.3 (23.2--46.9) & 50 min each way & air conditioning on recirculation & A & \citep{shen2020}$^{\mathrm{e}}$ \\
T2 & Coach, 49 seats, windows closed (Hunan) & 7/48 (or 7/46; 45 with seat positions) & 14.6 (6.1--27.8) & 2.5 h (200 min in the second report) & 1.72 l/s per person, measured & A & \citep{luo2020,ou2022} \\
T3 & Minibus, 18 seats (Hunan) & 2/12 (or 2/17) & 16.7 (2.1--48.4) & 60 min & 3.22 l/s per person, measured & B & \citep{luo2020,ou2022} \\
T4 & High-speed trains, pooled contact-tracing cohort (China) & 234/72{,}093 & 0.32 (0.28--0.37) & mean 2.1 h & not reported & A$^{*}$ & \citep{hu2021}$^{\mathrm{e}}$ \\
T5 & Long-haul flight, business cabin (London--Hanoi) & 12/20 & 60.0 (36.1--80.9) & 10 h & aircraft cabin; not quantified & A & \citep{khanh2020} \\
T6 & Airliner held on the ground, influenza (Alaska) & 72\,\% of passengers; 54 aboard (counts not transcribed) & -- & 3 h & none (system inoperative) & C & \citep{moser1979}$^{\mathrm{a}}$ \\
T7 & National cluster surveillance; commuter rail as a venue (Japan) & no denominator & -- & -- & -- & C & \citep{furuse2020} \\
\addlinespace[2pt]\multicolumn{8}{@{}l}{\emph{Restaurants}}\\
R1 & Restaurant, 18 tables, 89 patrons (Guangzhou) & 5/79$^{\ddagger}$ & 6.3 (2.1--14.2) & 82 min (index table) & 0.56--0.77 air changes per hour, measured & A & \citep{li2021,lu2020} \\
R2 & Restaurant, 9.2 m $\times$ 10.5 m (Jeonju) & 2/13 & 15.4 (1.9--45.4) & overlaps of 5 and 21 min & no ventilation system; air-conditioner jet & A & \citep{kwon2020} \\
\addlinespace[2pt]\multicolumn{8}{@{}l}{\emph{Choir and church}}\\
H1 & Choir rehearsal in a church hall (Skagit County) & 52/60 (32/60 confirmed) & 86.7 (75.4--94.1) & 2.5 h & 0.3--1.0 h$^{-1}$, estimated & A$^{*}$ & \citep{miller2021}$^{\mathrm{m}}$, \citep{hamner2020}$^{\mathrm{e}}$ \\
H2 & Church services, singer in a choir loft (Sydney) & 12/508 & 2.4 (1.2--4.1) & 1 h per service & minimal: fans off, doors and windows largely closed & B & \citep{katelaris2021} \\
\addlinespace[2pt]\multicolumn{8}{@{}l}{\emph{Outside the scope of the model (recorded for reference)}}\\
X1 & Cruise ship under quarantine (Yokohama) & 634/3{,}711$^{\dagger}$ & 17.1 (15.9--18.3) & weeks & -- & X & \citep{mizumoto2020} \\
X2 & Migrant-worker dormitories (Singapore) & 17{,}758/295{,}000$^{\dagger}$ & 6.0 (5.9--6.1) & weeks to months & -- & X & \citep{koh2020} \\
% <<< GENERATED records
\end{longtable}
}

\begin{figure}[tbp]
\centering
\includegraphics[width=\textwidth]{appB_dataset_attack_en}
\caption{Attack rates of the 23 records that have an exposed denominator, with exact (Clopper--Pearson)
95\,\% intervals, in the order of Table~\ref{appB_dataset:tab:records}. Filled circle: the count used in this
article; open circle: the alternative count (O2: 10 instead of 8 secondary cases; S2: 7 instead of 6;
T2: 7/46; T3: 2/17; H1: the 32 laboratory-confirmed cases only). $\dagger$: counts include the index or
primary case(s). `Multi-day' means that no duration in hours is recorded. A$*$: graded A without meeting
every grade-A criterion. The events differ in pathogen, variant, mask use, testing and number of
generations, so the attack rates are not comparable as levels of risk.}
% src: code/appB_dataset_fig_attack.py; data/appB_dataset_fig_attack.json (rows drawn)
\label{appB_dataset:fig:attack}
\end{figure}

\paragraph{Grades.}
Grade A (8 records) was defined as: a single index case, an enumerated exposed group, positions at seat or
zone level, a reported duration, and a geometry that is reported or measurable. Grade B (9 records) requires a
denominator and a duration with partial spatial information; grade C (8 records) covers aggregate,
multi-generation, abstract-only or incompletely enumerated events; the 2 out-of-scope records are graded X.
% src: data/bsc_outbreaks/outbreaks.csv (quality); data/appB_dataset_numbers.json -> grades
The re-check found that only 5 of the 8 grade-A records (T1, T2, T5, R1, R2) meet all the
grade-A criteria. The classroom C1 has no exposure duration in hours and no room size; the train cohort T4
pools 2,334 index patients and its source prints no cell denominators; and no seat or zone positions were
published for the choir H1. These three are marked A$*$. They remain in use (T4 is the only matrix of attack
rates by seat offset, H1 the benchmark for the level of the attack rate), but the dataset contains five, not
eight, single-index events with positions, a duration and a geometry.
% src: notes/VERIFICATION.md -> Section 4.4, outbreak dataset, item 1 (partially confirmed);
%      data/appB_dataset_numbers.json -> grade_A

\paragraph{Reading access.}
All 31 DOIs were confirmed in the DOI registry, and the bibliographic fields were taken from Crossref; the
re-check repeated both checks.
% src: data/bsc_outbreaks/tables/doi_verification.json; code/bsc_outbreaks/verify/results/source_check.json
The numbers were transcribed from the full text for 21 sources and from the accepted manuscript for one
\citep{miller2021}. Three sources \citep{shen2020,hu2021,hamner2020} were available only as an automated text
extraction of the publisher's web page, not as the PDF. Their numbers were compared with the abstracts and
extracted a second time in the re-check, with the same values, and two internal checks hold (the risk ratios recomputed from the zone counts of the bus,
Table~\ref{appB_dataset:tab:strata}, equal those printed by \citet{shen2020}; the train table is discussed
below). These three sources have not been checked against the PDF files. Six sources were read as abstracts
only \citep{nicholls2024,nardell1991,riley1978,moser1979,lan2021,jiang2021deptstore}; the five records that
rest on abstracts alone (O3, O5, S4, C4, T6) are graded C.
% src: data/appB_dataset_numbers.json -> sources (full 21, manuscript 1, extraction 3, abstract 6), records_resting_on_abstracts_only;
%      data/bsc_outbreaks/sources.md (Access); notes/VERIFICATION.md (Section 4.4, outbreak dataset)

\paragraph{Conflicting counts and an upper bound.}
Where sources disagree, both versions are recorded. The two investigations
of the Hunan coach give 7 cases among 48 passengers over 2.5\,h \citep{luo2020} and 7 among 46 over 200\,min
\citep{ou2022}, and they place the infected passengers differently; the minibus of the same outbreak has
2/12 against 2/17; the department store has 6 or 7 staff cases; the Zurich team has 8 or 10 first-generation
cases among 12; and 20 of the 52 choir cases are probable, not laboratory-confirmed.
% src: data/bsc_outbreaks/outbreaks.csv -> T2, T3, S2, O2, H1 (attack_rate_reported, quality_notes, n_infected_alt);
%      notes/VERIFICATION.md (Section 4.4, outbreak dataset: seat reconstructions of Luo and Ou differ)
For the Guangzhou restaurant the five secondary cases at the two neighbouring tables are an upper bound on
the number infected at the lunch. The first report of that outbreak states that at least one member of each
of the two families was infected in the restaurant and that the further cases in those families may be due
to transmission within the family \citep{lu2020}; the number infected at the lunch outside the index table
is therefore between 2 and 5. The counts marked $\ddagger$ in
Tables~\ref{appB_dataset:tab:records} and~\ref{appB_dataset:tab:strata} carry this caveat.
% src: notes/VERIFICATION.md -> Section 4.4, outbreak dataset (further points: Guangzhou counts are an upper bound);
%      notes/VERIFICATION.md (Section 4.4, outbreak dataset: Lu et al., onsets in Li et al. Table A.1)

\begin{table}[tbp]
\centering
\caption{Near/far contrasts available in the sources. Risk ratios with log-normal (Katz) 95\,\% intervals
and Fisher's exact test, computed from the counts shown. The two rows of T1 and the two rows of R1 are
alternative splits of the same people. $^{\mathrm{c}}$: one cell is zero and 0.5 was added to every cell.
$\ddagger$: upper bound (see text). $^{\mathrm{e}}$: pupils enrolled in grades 7--9 and 10--12 (581 and 583);
1,161 were tested, 580 of them in grades 10--12. For T5 the source defines the near stratum as
`$\le2$ seats away'. O1 and C2 are multi-generation outbreaks, and the position of the index
case in O1 is unknown.}
% src: data/bsc_outbreaks/outbreaks.csv -> C2 (attack_rate_reported: grades 197 + 197 + 187 = 581, 200 + 195 + 188 = 583; n_exposed 1161); data/checks/outbreak_dataset.json: jerusalem_senior_wing (583 persons, 580 tested; notes/VERIFICATION.md -> Section 4.4, outbreak dataset, further point 9); data/misc_checks.json (jerusalem; flight_label: underlined "<" before "2 seats away" in the stored copy of the source, which is not redistributed)
\label{appB_dataset:tab:strata}
\footnotesize
\setlength{\tabcolsep}{4pt}
\begin{tabular}{@{}l>{\raggedright}p{4.6cm}llll@{}}
\toprule
ID & Contrast (near against far) & Near: cases/$n$ (\%) & Far: cases/$n$ (\%) & Risk ratio (95\,\% CI) & $p$ \\
\midrule
% src: data/bsc_outbreaks/tables/near_far_risk_ratios.csv; recomputed and checked in code/appB_dataset_table.py
% >>> GENERATED strata (code/appB_dataset_table.py; do not edit by hand)
O1 & north wing against south wing (mapped seats) & 79/137 (57.7) & 4/62 (6.5) & 8.9 (3.4--23.3) & $5.2\times10^{-13}$ \\
C1 & front two rows against back three rows & 8/10 (80.0) & 4/14 (28.6) & 2.8 (1.2--6.8) & 0.036 \\
T1 & rows 5--11 (within about 2 m of the index) against rows 1--4 and 12--15 & 14/33 (42.4) & 9/34 (26.5) & 1.6 (0.8--3.2) & 0.20 \\
T1 & rows 6--10 against rows 1--5 and 11--15 (alternative split) & 11/23 (47.8) & 12/44 (27.3) & 1.8 (0.9--3.3) & 0.11 \\
T2 & rear rows 8--13 (index zone) against front rows 1--7 & 3/19 (15.8) & 4/26 (15.4) & 1.0 (0.3--4.1) & 1.0 \\
T5 & within 2 seats of the index against farther seats, business cabin & 11/12 (91.7) & 1/8 (12.5) & 7.3 (1.2--46.2) & $7.7\times10^{-4}$ \\
R1 & neighbouring tables against remote tables & 5/16$^{\ddagger}$ (31.2) & 0/63 (0.0) & 41.4$^{\mathrm{c}}$ (2.4--713) & $1.9\times10^{-4}$ \\
R1 & same air-conditioning zone against other zones (alternative split) & 5/11$^{\ddagger}$ (45.5) & 0/68 (0.0) & 63.2$^{\mathrm{c}}$ (3.7--1{,}071) & $2.0\times10^{-5}$ \\
C2 & grades 7--9 against grades 10--12 (separate wings) & 135/581 (23.2) & 18/583$^{\mathrm{e}}$ (3.1) & 7.5 (4.7--12.1) & $1.3\times10^{-26}$ \\
S1 & staff against customers & 11/120 (9.2) & 0/8{,}224 (0.0) & 1{,}563$^{\mathrm{c}}$ (92.7--26{,}381) & $3.4\times10^{-21}$ \\
% <<< GENERATED strata
\bottomrule
\end{tabular}
\end{table}

\begin{table}[tbp]
\centering
\caption{Negative controls: groups with a documented opportunity for exposure and almost no infection. The
last numerical column is the one-sided exact 95\,\% upper bound on the attack rate. The two restaurant rows
are alternative groupings of the same patrons. Access marks as in Table~\ref{appB_dataset:tab:records}.}
\label{appB_dataset:tab:controls}
\footnotesize
\begin{tabular}{@{}llllc@{}}
\toprule
Group & Record & Cases/exposed & Upper bound, \% & Source \\
\midrule
% src: data/bsc_outbreaks/tables/analysis_results.json -> negative_controls; bounds recomputed in code/appB_dataset_table.py
% >>> GENERATED controls (code/appB_dataset_table.py; do not edit by hand)
Supermarket customers & S1 & 0/8{,}224 & 0.036 & \citep{tian2021} \\
Other bus travelling to the same event & T1 & 0/60 & 4.9 & \citep{shen2020}$^{\mathrm{e}}$ \\
Restaurant, remote tables & R1 & 0/63 & 4.6 & \citep{li2021} \\
Restaurant, other air-conditioning zones & R1 & 0/68 & 4.3 & \citep{li2021} \\
Flight, premium-economy cabin & T5 & 0/35 & 8.2 & \citep{khanh2020} \\
Tested school contacts, masked classrooms & C3 & 5/728 & 1.44 & \citep{hershow2021} \\
Call-centre building, floors 7--9 & O1 & 1/595 & 0.79 & \citep{park2020} \\
% <<< GENERATED controls
\bottomrule
\end{tabular}
\end{table}

\paragraph{Spatial contrasts and negative controls.}
Table~\ref{appB_dataset:tab:strata} gives every near/far contrast that the sources allow, and
Table~\ref{appB_dataset:tab:controls} the groups used as negative controls. Three cautions apply. First, the
wing contrast of the call centre refers to the 84 case seats shown in the published plan; if the ten cases
without a seat had all worked in the south wing the ratio would be 2.6 (95\,\% interval 1.6--4.1), not 8.9.
% src: data/appB_dataset_numbers.json -> callcentre.rr, callcentre.rr_all_unmapped_south;
%      code/bsc_outbreaks/verify/results/callcentre_check.json -> wing_contrast.sens_all10_south
Second, the four single-index events with one split each (C1, T1, T2, T5) are statistically compatible with
a common risk ratio (Cochran's $Q=3.77$ on 3 degrees of freedom, $p=0.29$), so sorting events into
`near-field' and `well-mixed' ones is a reading of point estimates; the coach, with no front-to-rear
gradient, has a side-to-side contrast of 6/24 against 1/20 (risk ratio 5.0, 95\,\% interval 0.7--38).
% src: data/appB_dataset_numbers.json -> heterogeneity (cochran_Q 3.772, p 0.287), coach_side_contrast;
%      code/bsc_outbreaks/verify/results/extra_checks.json -> heterogeneity; stats_check.json -> risk_ratios (T2side)
Third, the levels are heterogeneous: over the nine SARS-CoV-2 events with a single index case and a single
exposure period, the cumulative hazard per hour, $-\ln(1-\text{attack rate})/\text{duration}$, spans a factor 34
when the 20 probable choir cases are counted and a factor 20 with confirmed cases only.
% src: data/bsc_outbreaks/tables/analysis_results.json -> hazard_per_hour_range (ratio 33.7);
%      code/bsc_outbreaks/verify/results/stats_check.json -> hazard_per_hour.confirmed_only_choir (ratio 19.97)

\paragraph{Digitisation of the call-centre floor plan.}
The seat map of record O1 was obtained from the floor plan published as Fig.~2 of \citet{park2020}; the
coordinates are derived from that figure and are not original data of this article. Desk rectangles were
segmented by colour and shape. The two open-plan wings contain 137 desks (north) and 62 desks (south). The
plan marks 84 case seats: 79 in the north wing, 4 in the south wing and 1 in the offices on the east side,
whose other desks were not digitised. The text of the source reports 94 cases among the 216 employees of
the floor, so 10 cases have no seat in the plan. The plan has no scale bar: coordinates are in units of the
desk pitch, their conversion to metres is unknown, and one employee per drawn desk is assumed.
% src: data/bsc_outbreaks/tables/park2020_callcentre_seats_digitized.csv; code/bsc_outbreaks/scripts/digitize_callcentre.py;
%      data/appB_dataset_numbers.json -> callcentre (north 79/137, south 4/62, case_seats_mapped 84, cases_unmapped 10)
The re-check typed the map again desk by desk from enlargements of the figure: the two readings match one to
one, no desk has a different case flag, and the largest difference in position is 3.8 pixels at a desk
pitch of 32.5 pixels. Figure~\ref{s8_outbreaks:fig:callcentre} is drawn from this file.
% src: code/bsc_outbreaks/verify/results/callcentre_check.json -> comparison_with_first_analysis (bijection true, 0 flags, 3.84 px);
%      data/appB_dataset_numbers.json -> callcentre.check_match, callcentre.desk_pitch_px

\paragraph{Reconstruction of the train cell counts.}
\citet{hu2021} print, for 23 seat offsets from the index passenger, an attack rate with a 95\,\% interval
but no counts. Two reconstructions were needed. On the publisher's page the five same-row values appear one
column to the left of where they belong (a contact cannot sit zero rows and zero columns away); with the
values moved to columns 1--5, the column means recomputed from the cells differ from the printed ones by
0.02 percentage points in total, against 1.17 for the layout as printed. Then, for each cell, the re-check
searched for the integer pair (cases, contacts) whose rate and Wilson interval reproduce the printed values;
we checked that every recovered pair does so to within 0.012 percentage points.
% src: data/bsc_outbreaks/tables/hu2021_train_attack_matrix.csv (realigned cells flagged);
%      code/bsc_outbreaks/verify/results/train_matrix_reconstruction.csv; verify/scripts/v02_stats.py;
%      data/appB_dataset_numbers.json -> train (column_mean_error_pp_realigned 0.024, _as_printed 1.170, max_abs_error_pp_vs_printed 0.0112)
The recovered cells sum to 230 cases among 72,692 contacts, against the 234 among 72,093 reported for the
cohort; cell denominators range from 1,033 to 5,420, and the 11 cells with at most three cases are the
least certain. The calibration of Section~\ref{M-s8_outbreaks:sec} uses these reconstructed counts. The source
itself warns that relatives and friends travelling together may have infected each other before or after the
journey, which inflates the adjacent-seat cell (3.53\,\%) most.
% src: data/appB_dataset_numbers.json -> train (230, 72692, 234, 72093, cell_n_min 1033, cell_n_max 5420, cells_with_at_most_3_cases 11);
%      code/bsc_validation/PREREGISTRATION.md, amendment V5; notes/VERIFICATION.md (Section 4.4, outbreak dataset, further point 3: Hu et al., companions)

\paragraph{Limitations and release of the first dataset.}
Investigated outbreaks are mostly large ones, so the attack rates over-represent highly infectious index
cases; records C3 and T4 are closer to ordinary exposures. Variants, mask use, vaccination and testing
differ between records and are not harmonised. Positions are zonal in most records: only T4 gives a matrix
of offsets and only O1 is resolved to seats.
% src: data/bsc_outbreaks/outbreaks.csv (quality_notes)
The released repository contains the dataset (\texttt{outbreaks.csv}, \texttt{outbreaks.json}), the list of
sources with the access level of each, the auxiliary tables (strata, risk ratios, the train matrix and its
reconstruction, the digitised seat coordinates) and the scripts that build and check them.
Copies of the source texts and figures used for transcription are not redistributed.

\subsection{The second dataset: eleven records of ten outbreaks with positions}\label{appB_dataset:sec:second}

\begin{table}[tbp]
  \centering\footnotesize
  \setlength{\tabcolsep}{2.5pt}
  \caption{The eleven records of the second dataset (Section~\ref{M-s8_outbreaks:sec:second}); W1 and W2 are the same
  ward with the same index patient. Near: within two
  rows of a source (aircraft), in the index patient's cubicle or bay (ward), within 8\,m (processing line).
  Risk ratio near/far with its 95\,\% interval. Grade A: identified source or sources, enumerated exposed
  group, positions for the whole group; B: more than one generation of cases. Role: fixed by the pathogen
  before any model was run.}
  \label{appB_dataset:tab:second}
  % generated: data/v2_tab_events2.tex (code/v2_numbers.py <- data/bsc_validation2/a1_more_outbreaks/07_descriptive.json; code/bsc_validation2/a1_more_outbreaks/PREREGISTRATION.md section 2.1)
  \begin{tabular}{@{}l>{\raggedright\arraybackslash}p{2.65cm}>{\raggedright\arraybackslash}p{1.75cm}>{\raggedright\arraybackslash}p{2.05cm}>{\raggedright\arraybackslash}p{2.2cm}>{\raggedright\arraybackslash}p{1.5cm}>{\raggedright\arraybackslash}p{1.85cm}c>{\raggedright\arraybackslash}p{1.55cm}@{}}
    \toprule
    & Event & Pathogen & Setting, duration & Position data & Near vs far & Risk ratio & Gr. & Role \\
    \midrule
    \tabinput{data/v2_tab_events2}
    \bottomrule
  \end{tabular}

  \smallskip
  \begin{minipage}{0.96\textwidth}\footnotesize
  $^{a}$\,With a continuity correction of 0.5 (a stratum without a case).
  \end{minipage}
\end{table}

\paragraph{Search and inclusion.}
The second dataset was assembled after the first test, by queries to Europe PMC for outbreaks with a seat map or
with counts by distance for the whole exposed group, and from the references of \citet{hertzberg2016risk}.
Eleven events were kept (Table~\ref{appB_dataset:tab:second}); the meat-processing line is record O4 of the
first dataset, now with its counts by distance. The events found and not used, with the reason, were fixed in
the protocol before any model was run: no identified source on board, one or two cases only, exposure before
the flight, denominators missing by seat, or positions not published (among them the aircraft outbreak of
\citet{moser1979} and the measles outbreak analysed by \citet{riley1978}). No outbreak in a
courtroom, a fitness class or a restaurant with positions for the whole exposed group was found in openly
accessible form. The search is not a systematic review, and the re-check did not repeat it.
% src: data/outbreaks2/SOURCES.md ("Events catalogued but not modelled"); code/bsc_validation2/a1_more_outbreaks/PREREGISTRATION.md section 2.2; notes/VERIFICATION.md (Section 4.6, second outbreak test, item 16: unverifiable)

\paragraph{Transcription and its checks.}
Seat maps were transcribed by eye from the published figures, and each map was checked against the totals
printed in its source before any model was run. The totals agree exactly for the flights from Hong Kong, from
Los Angeles and from Tel Aviv. They differ by one to four persons, in far strata, for three flights: 216
passengers at risk with a seat against 220 in the text for the flight from Cancun; 141 persons on the map
against 142 for the flight to Naha; 110 passengers of the mid cabin against 111 or 112 for the flight to
Perth. Two full texts are not openly accessible: the flight from Hong Kong uses the seat map as redrawn by
\citet{hertzberg2016risk}, and the tuberculosis flight only the two strata of the abstract, with seats placed
at random inside them. The counts of the ward and of the processing line were typed from the text and from an
appendix table of the sources.
% src: code/bsc_validation2/a1_more_outbreaks/POSTHOC_LOG.md (P3); data/outbreaks2/SOURCES.md
The re-check re-read every seat symbol of the four held-out flights and of the flight from Hong Kong
against the stored figures, re-typed the counts of the processing line, and fetched the abstract and
full-text numbers of the other events again; all stratum counts match. Two calibration maps (the flights from
Los Angeles and from Cancun) were not re-read seat by seat. Three limits of the data are not in the counts:
about 22 passengers of the flight to Naha who were lost to follow-up are counted as non-cases, mostly in far
rows; most non-cases of the flight from Tel Aviv were never tested; and the published figure of the ward
gives each student's bed, which the analysis replaced by random beds within the three strata (using the exact
beds changes the pooled score against the well-mixed model by $+0.13$ and nothing else). The near zone of the
flight to Naha is not the one its source analysed: the source's own group `within two rows' has 14 persons,
against the 23 in rows 21--25 used here. And one flight (AF171) was excluded as a secondary tabulation, although
two calibration events are themselves secondary in that sense (the flight from Hong Kong, from a redrawn map,
and the tuberculosis flight, from an abstract); with two cases it would change nothing.
% src: data/checks/outbreaks2.json: naha (about 22 lost to follow-up; 14 persons within two rows in the source, 23 in rows 21-25 here), ward_exact_beds (+0.13), AF171 (2 cases); notes/VERIFICATION.md (Section 4.6, second outbreak test, item 2: confirmed, with the caveats listed; further point 7; item 16: AF171 dropped as secondary tabulation, CA112 and the tuberculosis flight from a redraw and an abstract)
The released tables give, for every event, the strata, the persons at risk and the cases, and for the seat-map
events the transcribed position of every person, with the reference and the figure or table of origin.

